\documentclass[conference]{IEEEtran}
\IEEEoverridecommandlockouts
% The preceding line is only needed to identify funding in the first footnote. If that is unneeded, please comment it out.
%Template version as of 6/27/2024

\usepackage{cite}
\usepackage{amsmath,amssymb,amsfonts}
\usepackage{algorithmic}
\usepackage{graphicx}
\usepackage{textcomp}
\usepackage{xcolor}
\usepackage{fontspec}
\usepackage[hidelinks]{hyperref}
\def\BibTeX{{\rm B\kern-.05em{\sc i\kern-.025em b}\kern-.08em
    T\kern-.1667em\lower.7ex\hbox{E}\kern-.125emX}}

\setmainfont{Times New Roman}
\setsansfont{Arial}
\setmonofont{Courier New}

\begin{document}

\title{Operating-State-Constrained Weather-Power Fusion Network for Regional Wind Farm Cluster Forecasting
\thanks{This research is supported by Smart-Grid National Science and Technology Major Project (Grant No. 2025ZD0805500).}}

\author{\IEEEauthorblockN{1\textsuperscript{st} Ao Luo}
\IEEEauthorblockA{\textit{School of Information Science and}\\
\textit{Technology}\\
\textit{University of Science and}\\
\textit{Technology of China}\\
Hefei 230022, China\\
aoluo@mail.ustc.edu.cn}
\and
\IEEEauthorblockN{2\textsuperscript{nd} Hongwei Zhao\textsuperscript{*}}
\IEEEauthorblockA{\textit{School of Information Science and}\\
\textit{Technology}\\
\textit{University of Science and}\\
\textit{Technology of China}\\
Hefei 230022, China\\
hwzhao@ustc.edu.cn}
\and
\IEEEauthorblockN{3\textsuperscript{rd} Lei Liu}
\IEEEauthorblockA{\textit{School of Information Science and}\\
\textit{Technology}\\
\textit{University of Science and}\\
\textit{Technology of China}\\
Hefei 230022, China\\
liulei13@ustc.edu.cn}
\and[\hfill\break\mbox{}\hfill]
\IEEEauthorblockN{4\textsuperscript{th} Qiuju Chen}
\IEEEauthorblockA{\textit{School of Information Science and}\\
\textit{Technology}\\
\textit{University of Science and}\\
\textit{Technology of China}\\
Hefei 230022, China}
\and
\IEEEauthorblockN{5\textsuperscript{th} Bin Li}
\IEEEauthorblockA{\textit{School of Information Science and}\\
\textit{Technology}\\
\textit{University of Science and}\\
\textit{Technology of China}\\
Hefei 230022, China}
}

\maketitle

\begin{abstract}
Intraday regional wind farm cluster power forecasting is an important basis for power dispatch, reserve allocation, and renewable energy integration. Weather-driven forecasting models usually rely on future weather forecasts to describe wind-field evolution during the forecast day. However, under the same weather situation, actual power output is also affected by the station-level generation level, short-term inertia, and synchronized multi-site variations before the forecast origin; if a model only focuses on future weather, it can produce deviations inconsistent with the true operating state. This paper proposes an operating-state-constrained weather-power fusion network. The method extracts weather-driven representations from meteorological forecasts, encodes station-level measured power available before the forecast origin as an operating-state representation, and uses this representation to constrain weather-information reading during decoding. Experiments on three years of data from 58 wind farms in Jilin Province show that the proposed method achieves a normalized MAE of 0.0779, a normalized RMSE of 0.1065, and an \(R^2\) of 0.8036 for intraday aggregate power forecasting. Compared with Weather-only, the MAE and RMSE are reduced by 6.84\% and 6.57\%, respectively; both errors are reduced by approximately 5\% relative to the fusion baselines. Horizon-wise analysis shows that the gain is most pronounced at the beginning of the forecast day, with first-horizon MAE reduced by 47.94\%. These results indicate that true power states before the forecast origin can effectively constrain the weather-to-power mapping and improve state adaptability in regional wind farm cluster forecasting.
\end{abstract}

\begin{IEEEkeywords}
regional wind farm cluster forecasting, wind power forecasting, weather-power fusion, operating-state constraints, multi-step forecasting
\end{IEEEkeywords}

\section{Introduction}

Intraday regional wind farm cluster power forecasting is an important basis for power dispatch, reserve allocation, and renewable energy integration. This study takes aggregate power from multiple wind farms in the region as the prediction target and investigates its evolution over the forecast day. Future meteorological forecasts describe wind-field evolution over the forecast period, whereas measured power before the forecast origin records the current cluster operating state; the two information sources therefore play different roles in forecasting. Regional weather propagation changes the wind-energy input across farms and thereby drives intraday aggregate-power variations. Meanwhile, the current generation level, short-term inertia, and synchronized multi-site variations affect the initial level and near-term evolution of the power trajectory; recovery from curtailment may further lead to different actual outputs under similar weather conditions. Accordingly, the model needs to account for both future wind-field changes and the operating state at the forecast origin.

Gridded weather forecasts are available throughout the forecast day, whereas measured power is available only before the forecast origin. The weather forecast provides the external drivers of subsequent power changes but does not contain the actual power state at the forecast origin; this state must be characterized from pre-origin power records. The current generation level, short-term inertia, and synchronized multi-site variations reflected in recent measured power can therefore serve as an operating-state constraint when reading future-weather information. This paper addresses how to introduce this constraint into future-weather-driven forecasting.

Existing wind power forecasting studies have improved prediction performance through wind-speed sequence modeling \cite{ref1} and gridded spatial representations \cite{ref2}. Cluster-oriented studies have developed adaptive spatiotemporal graph models \cite{ref3} and hierarchical graph models \cite{ref4}, as well as pattern-aware graph predictors \cite{ref5} and global-information graph encoders \cite{ref6}. Dynamic graph updates \cite{ref7} and learnable graph structures \cite{ref8} have also been investigated. Other directions include trend-aware graph switching \cite{ref9}, multi-model fusion \cite{ref10}, trend-based spatial integration \cite{ref11}, and continual learning for changing graph inputs \cite{ref12}. For future-weather information, multimodal Transformer fusion \cite{ref13} and spatiotemporal attention fusion \cite{ref14} have been developed; other work addresses NWP temporal lag \cite{ref15}, wind-speed offsets \cite{ref16}, predictive uncertainty \cite{ref_liu_bayesian_uncertainty}, multi-source spatiotemporal fusion \cite{ref17}, or spatial correlation with corrected NWP \cite{ref18}. The focus of this paper is to explicitly use station-level power states observable before the forecast origin in future-weather-driven intraday forecasting. Recent power records contain the current generation level, short-term inertia, and synchronized multi-site variations, and can serve as operating-state constraints for future-weather reading, thereby improving forecasting accuracy.

To impose this constraint, forecasting is organized as weather-context construction, forecast-origin state encoding, and state-conditioned decoding. Recent station power becomes dynamic context that conditions how queries read weather memory.

The main contributions are as follows:
\begin{itemize}
\item Recent measured power is formulated as a forecast-origin operating-state constraint, distinct from future weather as the main driver of intraday evolution.
\item Memory-level conditioning converts multi-scale station power into context tokens that participate directly in weather-memory reading.
\item Weather-only and fusion comparisons, true/random/zero-power controls, and horizon-wise analyses on 58 wind farms jointly verify the value and time-varying influence of the true origin state.
\end{itemize}

\begin{figure*}[htbp]
\centerline{\includegraphics[width=0.78\textwidth]{../../paper_image_use/method_new/exported/Fig01.pdf}}
\caption{Business time protocol. Recent station-level power is only available before the forecast day starts; the meteorological forecast is divided into a historical segment and a target segment according to its temporal relation to the forecast day.}
\label{fig:time-protocol}
\end{figure*}

\section{Related Work and Difference}

Studies on regional or cluster wind power forecasting have investigated temporal variations and spatial correlations among wind farms through wind-speed sequence modeling, regular-grid embedding, and graph-structure modeling. Early deep learning methods used CNN-LSTM to model wind-speed temporal features \cite{ref1}, while other work embedded turbines into regular grids so that convolutional networks could process spatial distributions \cite{ref2}. For wind farm clusters, Wang et al. adaptively adjusted the graph topology during training and combined spatiotemporal graph convolution, improved multi-head attention, and multi-task outputs \cite{ref3}. Li et al. partitioned wind farms into hierarchical subregions and combined graph modeling with extreme-value-aware correction \cite{ref4}. Liu et al. identified power-fluctuation patterns and trained pattern-specific spatiotemporal graph predictors \cite{ref5}. Yang et al. constructed multiple global-information graphs and switched among adaptive graph convolution models according to dynamic correlation states \cite{ref6}. A dynamic graph neural network was further used to update spatial correlations through node embeddings \cite{ref7}, while a graph deviation attention network learned graph structures and dynamically corrected errors during load peak and valley periods \cite{ref8}. Fluctuation-correlation graphs with an NWP trend-aware switching mechanism were introduced for deterministic and probabilistic cluster forecasting \cite{ref9}. D-stacking combined spatiotemporal graph features with dual multi-model fusion for rolling multi-step prediction \cite{ref10}. Trend categorization and spatial information integration were used to organize cluster-level day-ahead inputs \cite{ref11}, and a globally aware dynamic spatiotemporal graph continual-learning framework was proposed to adapt to changing graph inputs without forgetting historical dynamics \cite{ref12}. These studies demonstrate the value of multi-station power dynamics and spatial coupling. However, they mainly emphasize historical-power representation or cluster graph construction; how the true power state at the forecast origin should constrain future-weather reading remains less explicit.

For future-weather-driven multi-step wind power forecasting, M2TNet jointly learned multimodal representations and auxiliary tasks through a multi-task Transformer \cite{ref13}, while M2STAN extended multimodal fusion with spatiotemporal attention for multi-location prediction \cite{ref14}. Series-wise attention was introduced to handle the temporal lag of NWP variables \cite{ref15}, and an error-tolerability strategy was designed for NWP wind-speed offsets \cite{ref16}. Deep Bayesian modeling was used to quantify predictive uncertainty \cite{ref_liu_bayesian_uncertainty}. Multi-source spatiotemporal features were fused for cluster forecasting \cite{ref17}, and spatial correlation was combined with corrected NWP \cite{ref18}. For gridded meteorological fields, 3D CNNs directly extracted features from NWP data cubes \cite{ref19}; WindDragon provided an automated regional deep-learning framework \cite{ref20}; and deep CNN encoding was combined with conformalized regression forests for regional probabilistic forecasts \cite{ref21}. This paper uses Pangu-Weather as the future weather source \cite{ref22} and treats it as the main driving information for intraday power changes.

Informer introduced efficient attention for long-sequence time-series forecasting \cite{ref23}. PatchTST represented temporal sequences as patches \cite{ref24}, whereas Crossformer modeled cross-dimension dependencies in multivariate series \cite{ref25}. Non-stationary Transformers explicitly addressed distributional non-stationarity \cite{ref26}. In wind power forecasting, a feature-temporal Transformer provided interpretable multivariate modeling \cite{ref_liu_feature_temporal}, and TDFormer decoupled multiple dependencies for accurate multi-step prediction \cite{ref_liu_tdformer}. These methods indicate that attention mechanisms are suitable for long-range dependencies in multi-step forecasting. In operational wind power forecasting, the distinct roles of future weather drivers and the forecast-origin operating state still need to be formalized more explicitly.

Existing multi-source fusion studies usually emphasize feature complementarity among weather, power, and spatial information, with limited attention to how the power state at the forecast origin constrains future-weather reading. This paper generates dynamic context representations from station-level measured power before the forecast origin and introduces them into weather-memory reading, so that weather-information extraction at each forecast step is constrained by the current operating state.

\section{Problem Definition and Business-Time Alignment}

Let the forecast day be \(D\), with an hourly data resolution. To distinguish weather-product initialization from the power-forecast origin, let \(t_{\mathrm{init}}\) denote the Pangu-Weather initialization time, 08:00 Beijing time on day \(D-1\), and let \(t_0\) denote the power-forecast origin, 23:00 on day \(D-1\). We use \(L=16\) for the number of hourly observations in the power-state window, which spans \(t_0-L+1\) through \(t_0\), i.e., 08:00--23:00 on day \(D-1\). We use \(H=24\) for the number of intraday forecast steps. Forecast step \(h\) corresponds to \(t_0+h\), for \(h=1,\ldots,H\), so the target window spans 00:00--23:00 on day \(D\).

Let \(N\) be the number of wind farms. The pre-origin state context is \(\mathbf{P}^{c}\in\mathbb{R}^{N\times L}\). The regional aggregate-power target is \(\mathbf{y}=[y_{t_0+1},\ldots,y_{t_0+H}]\in\mathbb{R}^{H}\), and the model output is \(\hat{\mathbf{y}}\in\mathbb{R}^{H}\). The station adjacency matrix \(\mathbf{A}\in\mathbb{R}^{N\times N}\) is estimated from the same \(L\)-step power-state window, retaining positive Pearson correlations greater than 0.7 as edge weights.

Figure~\ref{fig:time-protocol} shows the business-time protocol. Within the single Pangu-Weather product initialized at \(t_{\mathrm{init}}\), lead times 0 through \(L-1\) are valid from 08:00 through 23:00 on day \(D-1\) and form the historical forecast segment \(\mathbf{X}^{h}\). Lead times \(L\) through \(L+H-1\) are valid over the target window on day \(D\) and form the target forecast segment \(\mathbf{X}^{f}\). The terms ``historical'' and ``target'' indicate whether a frame's valid time is before or within the target window; both segments belong to the same weather forecast generated at \(t_{\mathrm{init}}\). They are concatenated as \(\mathbf{X}=[\mathbf{X}^{h};\mathbf{X}^{f}]\). Thus, the power forecast issued at origin \(t_0\) uses only measured power through \(t_0\) and the previously initialized weather product, with no measured power from the target window.

Hourly sliding windows are used for training while preserving this relative alignment between inputs and targets. Validation and testing use strict daily samples following the fixed business-time protocol for day \(D\).

\begin{figure*}[htbp]
\centerline{\includegraphics[width=0.72\textwidth]{../../paper_image_use/method_new/exported/Fig02.pdf}}
\caption{Overall framework. Future weather forms the main memory, while recent station-level power generates dynamic context tokens and imposes an operating-state constraint on weather reading at the memory level.}
\label{fig:framework}
\end{figure*}

\section{Operating-State-Constrained Weather-Power Fusion Network}

\subsection{Overall Framework}

As shown in Fig.~\ref{fig:framework}, the network implements three functional stages. The Weather Encoder converts meteorological forecasts into queryable context carrying the main intraday driver. The Power-State Encoder summarizes only pre-origin station power into a sample-specific state. The Query Decoder then predicts all horizons in parallel by reading weather context conditioned on that state.

The model applies memory-level conditioning by mapping the encoded state into dynamic tokens in the weather-memory latent space, allowing each horizon query to read origin-state and future-weather information from the same fused memory. The following subsections introduce the weather-memory construction, power-state encoding, and memory-level fusion and query-decoding structures.

\subsection{Weather Memory Construction}

The weather memory construction process is shown in Fig.~\ref{fig:weather-memory}. The Weather Encoder takes \(\mathbf{X}=[\mathbf{X}^{h};\mathbf{X}^{f}]\in\mathbb{R}^{T\times C\times H_g\times W_g}\), where \(T=L+H\), and \(L\) and \(H\) are the power-state window length and intraday forecast length defined in Section~III, respectively. The historical segment \(\mathbf{X}^{h}\) and target segment \(\mathbf{X}^{f}\) contain \(L\) and \(H\) frames, respectively. The \(C=17\) original channels are divided into nine wind-field and eight background meteorological channels, and first-order lead-time differences are additionally computed for all channels. Separate CNN stems extract frame-wise spatial features from the three information groups:
\begin{equation}
\mathbf{F}^{u}_{t}=\phi_{u}(\mathbf{X}^{u}_{t}),\quad
\mathbf{F}^{b}_{t}=\phi_{b}(\mathbf{X}^{b}_{t}),\quad
\mathbf{F}^{\Delta}_{t}=\phi_{\Delta}(\Delta\mathbf{X}_{t}),
\end{equation}
Here, \(\mathbf{X}^{u}_{t}\) and \(\mathbf{X}^{b}_{t}\) select the wind-field and background meteorological channels of frame \(t\), respectively. The first-order lead-time difference is defined as
\begin{equation}
\Delta\mathbf{X}_{1}=\mathbf{0},\qquad
\Delta\mathbf{X}_{t}=\mathbf{X}_{t}-\mathbf{X}_{t-1},\quad t=2,\ldots,T.
\end{equation}
The frame-wise 2D convolutional stems \(\phi_{u}\), \(\phi_{b}\), and \(\phi_{\Delta}\) encode wind drivers, background meteorological conditions, and lead-time changes, respectively.

The three feature groups are fused by a \(1\times1\) block and then augmented with temporal and stage embeddings:
\begin{equation}
\mathbf{H}_{t}=\psi([\mathbf{F}^{u}_{t};\mathbf{F}^{b}_{t};\mathbf{F}^{\Delta}_{t}])
+\mathbf{e}^{time}_{t}+\mathbf{e}^{stage}_{t}.
\end{equation}
Here, \(\psi(\cdot)\) denotes the \(1\times1\) convolutional fusion block used to mix channel information from the three meteorological branches; \(\mathbf{H}_{t}\) is the fused hidden representation of the \(t\)-th meteorological frame; \(\mathbf{e}^{time}_{t}\) marks the relative position of the frame in the input sequence; and \(\mathbf{e}^{stage}_{t}\) distinguishes the historical and target segments of the meteorological forecast sequence.

The sequence \(\mathbf{H}_{1:T}\) is then fed into a spatiotemporal convolutional backbone \(\mathcal{E}_{w}\). The backbone consists of depthwise 3D spatiotemporal convolutions, spatial downsampling convolutions, and frame-wise spatial mixing modules. It models local spatiotemporal variations while preserving temporal length. The encoded features are first projected by a \(1\times1\times1\) memory projection and then pooled into a fixed-size spatiotemporal grid:
\begin{equation}
\mathbf{Z}^{w}=
\operatorname{Pool}\left(\operatorname{Proj}(\mathcal{E}_{w}(\mathbf{H}_{1:T}))\right)
\in\mathbb{R}^{T_m\times H_m\times W_m\times d}.
\end{equation}
Here, \(\operatorname{Proj}\) is the \(1\times1\times1\) convolutional projection, and \(\operatorname{Pool}\) denotes adaptive average pooling. To preserve the spatiotemporal positions of memory tokens, temporal, row, and column learnable position embeddings are added:
\begin{equation}
\mathbf{E}^{pos}_{\tau,i,j}=
\mathbf{p}^{time}_{\tau}+\mathbf{p}^{row}_{i}+\mathbf{p}^{col}_{j}.
\end{equation}
The final weather memory is defined as
\begin{equation}
\begin{aligned}
\mathbf{M}^{w}
&=\operatorname{Flatten}(\mathbf{Z}^{w}+\mathbf{E}^{pos})
\in\mathbb{R}^{L_w\times d},\\
L_w&=T_mH_mW_m.
\end{aligned}
\end{equation}
Here, Flatten converts the pooled grid into a token sequence, and
\(L_w\) is the number of weather-memory tokens. This memory preserves both forecast segments and local spatiotemporal structure, allowing queries at different forecast steps to read information from different lead times and spatial regions.

\begin{figure*}[htbp]
\centerline{\includegraphics[width=0.72\textwidth]{../../paper_image_use/method_new/exported/Fig03.pdf}}
\caption{Weather memory construction. Decoupled meteorological branches extract wind-field, background-field, and meteorological-change features, and then construct spatiotemporal weather memory that can be read by future queries.}
\label{fig:weather-memory}
\end{figure*}

\subsection{Power-State Encoding}

The power-state encoding process is shown in Fig.~\ref{fig:power-state}. The Power-State Encoder takes station-level power available before the forecast origin, \(\mathbf{P}^{c}\in\mathbb{R}^{N\times L}\), and the station adjacency matrix \(\mathbf{A}\in\mathbb{R}^{N\times N}\) as input, where \(N\) is the number of wind farms and \(L\) is the power-state window length defined in Section~III.

The model encodes station-level power at point, patch, and global temporal scales, respectively characterizing point-wise spatial states, local ramp patterns, and the overall generation trend within the context range. The local patches are defined by
\begin{equation}
\label{eq:patch-start}
a_j=1+(j-1)r,\qquad j=1,\ldots,J,
\end{equation}
\begin{equation}
\label{eq:patch-index}
\mathcal{I}_{j}=\{a_j,\ldots,a_j+s-1\},
\end{equation}
\begin{equation}
\label{eq:patch-count}
J=\left\lfloor\frac{L-s}{r}\right\rfloor+1.
\end{equation}
Here, \(s\) is the patch length, \(r\) is the sliding stride, \(a_j\) is the starting index of patch \(j\), \(\mathcal{I}_j\) is its temporal index set, and \(J\) is the number of valid patches. None of these patches uses target-period power. The three token types are written as
\begin{equation}
\label{eq:point-token}
\mathbf{o}^{pt}_{t}=g_{pt}(\operatorname{GCN}_{pt}(\mathbf{P}^{c}_{:,t},\mathbf{A})),
\quad t=1,\ldots,L,
\end{equation}
\begin{equation}
\label{eq:patch-token}
\mathbf{o}^{pa}_{j}
 =g_{pa}(\operatorname{GCN}_{pa}(\mathbf{P}^{c}_{:,\mathcal{I}_{j}},\mathbf{A})),
\quad j=1,\ldots,J,
\end{equation}
\begin{equation}
\label{eq:global-token}
\mathbf{o}^{gl}
 =g_{gl}(\operatorname{GCN}_{gl}(\mathbf{P}^{c},\mathbf{A})).
\end{equation}
The graph-convolution node outputs are flattened across farm and feature dimensions and then mapped to the \(d\)-dimensional token space by linear projections \(g_{pt}\), \(g_{pa}\), and \(g_{gl}\), respectively.

To integrate information across the three temporal scales, the power encoder adopts bottom-up and top-down cross-scale attention. After positional encodings are added, the three token groups are denoted as \(\mathbf{H}^{pt}\), \(\mathbf{H}^{pa}\), and \(\mathbf{H}^{gl}\). Cross-scale interaction first lets patch tokens read point tokens and the global token read patch tokens; then the global token modulates patch tokens, and patch tokens further modulate point tokens:
\begin{equation}
\mathbf{H}^{pa}_{bu}=\operatorname{CSA}_{pt\rightarrow pa}
(\mathbf{H}^{pa},\mathbf{H}^{pt}).
\end{equation}
\begin{equation}
\mathbf{H}^{gl}_{bu}=\operatorname{CSA}_{pa\rightarrow gl}
(\mathbf{H}^{gl},\mathbf{H}^{pa}_{bu}).
\end{equation}
\begin{equation}
\mathbf{H}^{pa}_{td}=\operatorname{CSA}_{gl\rightarrow pa}
(\mathbf{H}^{pa}_{bu},\mathbf{H}^{gl}_{bu}).
\end{equation}
\begin{equation}
\mathbf{H}^{pt}_{td}=\operatorname{CSA}_{pa\rightarrow pt}
(\mathbf{H}^{pt},\mathbf{H}^{pa}_{td}).
\end{equation}
Here, subscripts \(bu\) and \(td\) denote bottom-up and top-down interaction outputs, respectively; \(\operatorname{CSA}_{source\rightarrow target}\) denotes a cross-scale attention block that uses target tokens as queries and source tokens as keys and values, followed by residual connections, LayerNorm, and an FFN. The fused cross-scale token sequence is
\begin{equation}
\mathbf{S}^{p}=[\mathbf{H}^{gl}_{bu};\mathbf{H}^{pa}_{td};\mathbf{H}^{pt}_{td}]
\in\mathbb{R}^{(1+J+L)\times d}.
\end{equation}
After adding a fusion positional encoding, the sequence is passed through a Transformer encoder, flattened, and projected into a sample-level power-state vector:
\begin{equation}
\mathbf{z}^{p}=f_{p}(\operatorname{vec}(\operatorname{Transformer}(\mathbf{S}^{p})))
\in\mathbb{R}^{d}.
\end{equation}
Here, \(\operatorname{vec}\) flattens the token sequence, and \(f_p\) is the state projection head that maps it to \(d\) dimensions. The resulting sample-specific operating-state representation compresses the current generation level, local changes, and inter-station coupling for subsequent weather-memory reading.

\begin{figure*}[htbp]
\centerline{\includegraphics[width=0.66\textwidth]{../../paper_image_use/method_new/exported/Fig04.pdf}}
\caption{Power-state encoding. The multi-scale GCN extracts current operating state, local variations, and multi-site coupling from recent station-level power.}
\label{fig:power-state}
\end{figure*}

\subsection{Memory-Level Fusion and Query Decoding}

To write the power state at the memory level, a two-layer feed-forward power-to-context adapter \(a_{\theta}\) maps the power-state vector into \(K\) dynamic power context tokens:
\begin{equation}
\mathbf{C}^{p}=\operatorname{reshape}(a_{\theta}(\mathbf{z}^{p}))\in\mathbb{R}^{K\times d}.
\end{equation}
These tokens are jointly generated from the same power-state vector and have no one-to-one correspondence with historical power steps. They are then prepended to the weather memory:
\begin{equation}
\mathbf{M}^{pc}=[\mathbf{C}^{p};\mathbf{M}^{w}]\in\mathbb{R}^{(K+L_w)\times d}.
\end{equation}

This prepended injection allows the power state to participate in cross-attention at all forecast steps in the form of memory tokens. Query vectors for different forecast steps can adaptively read power context and weather memory from the same fused memory, allowing the power state to act as an operating-state constraint during parallel decoding.

The Query Decoder starts from learnable future queries
\(\mathbf{Q}^{(0)}\in\mathbb{R}^{H\times d}\), with one query for each forecast step. The decoder comprises stacked decoder blocks. For \(\ell=1,\ldots,L_d\), each block consists of query self-attention, query-to-memory cross-attention, and an FFN with residual connections and LayerNorm. The self-attention models dependencies among different forecast steps, while the cross-attention reads weather memory and power context from the fused memory \(\mathbf{M}^{pc}\):
\begin{equation}
\label{eq:decoder-self}
\tilde{\mathbf{Q}}^{(\ell)}=\operatorname{LN}\!\left(
\mathbf{Q}^{(\ell-1)}+\operatorname{MSA}(\mathbf{Q}^{(\ell-1)})\right).
\end{equation}
\begin{equation}
\label{eq:decoder-cross}
\bar{\mathbf{Q}}^{(\ell)}=\operatorname{LN}\!\left(
\tilde{\mathbf{Q}}^{(\ell)}+\operatorname{MCA}(\tilde{\mathbf{Q}}^{(\ell)},\mathbf{M}^{pc})\right).
\end{equation}
\begin{equation}
\label{eq:decoder-ffn}
\mathbf{Q}^{(\ell)}=\operatorname{LN}\!\left(\bar{\mathbf{Q}}^{(\ell)}
+\operatorname{FFN}(\bar{\mathbf{Q}}^{(\ell)})\right).
\end{equation}
Here, \(\operatorname{MSA}(\mathbf{Q})=\operatorname{MHA}(\mathbf{Q},\mathbf{Q},\mathbf{Q})\) denotes multi-head self-attention, while \(\operatorname{MCA}(\mathbf{Q},\mathbf{M})=\operatorname{MHA}(\mathbf{Q},\mathbf{M},\mathbf{M})\) denotes multi-head cross-attention with queries \(\mathbf{Q}\) and memory keys/values \(\mathbf{M}\). Equations~\eqref{eq:decoder-self}--\eqref{eq:decoder-ffn} respectively perform horizon-wise self-attention, conditioned-memory reading, and feed-forward updating; \(\tilde{\mathbf{Q}}^{(\ell)}\) and \(\bar{\mathbf{Q}}^{(\ell)}\) are the two intermediate query representations in decoder block \(\ell\).
The final prediction is
\begin{equation}
\hat{\mathbf{y}}=g_{\phi}(\operatorname{LN}(\mathbf{Q}^{(L_d)})).
\end{equation}
Here, \(g_{\phi}\) is the prediction head that maps each horizon query representation to the corresponding regional aggregate power output, and is implemented by a linear layer, GELU, dropout, and an output linear layer.
Training uses only the final regional aggregate power forecast loss:
\begin{equation}
\mathcal{L}=\lambda_1\,\operatorname{MAE}(\hat{\mathbf{y}},\mathbf{y})
+\lambda_2\,\operatorname{RMSE}(\hat{\mathbf{y}},\mathbf{y}).
\end{equation}
\(\lambda_1\) and \(\lambda_2\) balance the MAE and RMSE terms, and are set to 0.6 and 0.4, respectively.

All future queries are decoded in parallel, avoiding autoregressive error accumulation. The training objective directly matches the final intraday regional wind power forecasting task.

\section{Experiments}

\subsection{Data and Settings}

This paper constructs an experimental dataset from a regional wind farm cluster in Jilin Province. The power data cover 2022--2024, and the prediction target is the regional aggregate output of 58 wind farms. The meteorological input is the hourly Pangu-Weather gridded forecast covering Jilin and its surrounding area, with a grid size of \(19\times33\) and a spatial resolution of \(0.25^\circ\). The 17 input channels include 10 m wind fields, wind fields at different pressure levels, temperature, humidity, geopotential height, and mean sea-level pressure. All power and meteorological data are aligned in Beijing time.

The training, validation, and test sets are isolated by calendar year and use data from 2022, 2023, and 2024, respectively. Sliding-window samples are used for training, whereas strict daily samples are used for validation and testing. All models use the same data split, normalization, and intraday forecasting target. Target power is scaled by the maximum forecast-period power in the training set; historical station power is standardized per station, and meteorological inputs are standardized per channel. The validation and test sets reuse only training-set statistics. The evaluation metrics include normalized MAE and RMSE, engineering-scale MAE and RMSE, and \(R^2\); all reported results are averages over 12 independent runs. The proposed method uses the latest \(L\) steps of station-level power as the operating-state context, estimates the adjacency matrix from the same window, and predicts \(H\) intraday steps.

To examine the gain from operating-state constraints, this paper first sets Weather-only as a direct comparison model that shares the same weather input and query decoder. CNN, 3D CNN, ViT, Swin, and Video Transformer are also introduced as meteorological-encoder baselines to observe pure-weather forecasting performance under different gridded-weather encoding strategies.

To ensure consistent use of temporal information, all weather-only baselines take the complete meteorological forecast (historical segment) and meteorological forecast (target segment) as input. CNN, ViT, and Swin use frame-wise spatial encoding followed by a temporal Transformer, whereas the 3D CNN and Video Transformer perform direct spatiotemporal encoding. iTransformer-Fusion and PatchTST-Fusion add a historical-power sequence branch to Weather-only and combine the two outputs through a softmax-gated MLP, providing output-level fusion comparisons for the proposed memory-level state injection.

\begin{figure*}[t]
\centerline{\includegraphics[width=\textwidth]{figures_short/horizon_comparison_en.pdf}}
\caption{Horizon-wise MAE and gap. (a) MAE. (b) MAE gap between Weather-only and the proposed method; a positive value indicates lower error for the proposed method.}
\label{fig:horizon}
\end{figure*}

\subsection{Overall Forecasting Performance}

Table~\ref{tab:benchmark-short} reports the overall results on the independent 2024 test set. The proposed method achieves the lowest normalized MAE, the lowest normalized RMSE, and the highest \(R^2\), with values of 0.0779, 0.1065, and 0.8036, respectively. The corresponding engineering-scale MAE and RMSE are 383.5 MW and 524.2 MW. Weather-only is the lowest-error pure-weather model, indicating that complete future meteorological forecasts already provide the main predictive signal.

\begin{table}[htbp]
\caption{Overall forecasting performance}
\begin{center}
\footnotesize
\begin{tabular}{lccc}
\hline
Model & MAE \(\downarrow\) & RMSE \(\downarrow\) & \(R^2\) \(\uparrow\) \\
\hline
Ours & \textbf{0.0779} & \textbf{0.1065} & \textbf{0.8036} \\
iTrans.-Fusion & 0.0821 & 0.1121 & 0.7821 \\
PatchTST-Fusion & 0.0821 & 0.1120 & 0.7824 \\
Weather-only & 0.0836 & 0.1140 & 0.7749 \\
3D CNN & 0.0856 & 0.1172 & 0.7619 \\
CNN & 0.0916 & 0.1248 & 0.7302 \\
Video Trans. & 0.0937 & 0.1265 & 0.7228 \\
ViT & 0.0941 & 0.1278 & 0.7167 \\
Swin & 0.0947 & 0.1279 & 0.7162 \\
\hline
\end{tabular}
\label{tab:benchmark-short}
\end{center}
\end{table}

Weather-only has the lowest error among the pure-weather models, while both output-level fusion baselines remain weaker than the proposed method after incorporating historical power. Relative to Weather-only, the proposed method reduces MAE and RMSE by 6.84\% and 6.57\%, respectively. Relative to iTransformer-Fusion, the reductions are 5.16\% and 5.05\%; relative to PatchTST-Fusion, they are 5.19\% and 4.99\%. These results show that using true pre-origin power states during memory reading yields lower errors than fusing an independent historical-power branch at the output level.

\subsection{Horizon-Wise Fusion Effect}

Table~\ref{tab:horizon-short} and Fig.~\ref{fig:horizon} show the error differences between the proposed method and Weather-only at different forecast steps, where \(h\) denotes the forecast-step index. At \(h=1\), the proposed method achieves an MAE of 0.0445, lower than the 0.0856 of Weather-only; the corresponding RMSE values are 0.0626 and 0.1138.

\begin{table}[htbp]
\caption{Horizon-wise error comparison}
\begin{center}
\footnotesize
\begin{tabular}{llccccc}
\hline
Model & Metric & \(h=1\) & \(h=6\) & \(h=12\) & \(h=H\) & Avg \\
\hline
Ours & MAE & \textbf{0.0445} & \textbf{0.0759} & \textbf{0.0742} & \textbf{0.0955} & \textbf{0.0779} \\
Weather & MAE & 0.0856 & 0.0806 & 0.0754 & 0.0973 & 0.0836 \\
Ours & RMSE & \textbf{0.0626} & \textbf{0.0974} & \textbf{0.1050} & \textbf{0.1215} & \textbf{0.1055} \\
Weather & RMSE & 0.1138 & 0.1066 & 0.1082 & 0.1245 & 0.1138 \\
\hline
\end{tabular}
\label{tab:horizon-short}
\end{center}
\end{table}

The step-wise curves show that the benefit of operating-state constraints is concentrated near the beginning of the forecast day and reduces step-averaged MAE from 0.0836 to 0.0779. At \(h=1\), MAE and RMSE are reduced by 47.94\% and 45.00\% relative to Weather-only; at the final step \(h=H\), the reductions remain 1.80\% and 2.40\%. The share of query cross-attention assigned to power context decreases from 39.56\% at \(h=1\) to 6.56\% at \(h=H\), consistent with the larger near-end error reduction. The average RMSE in the table is the arithmetic mean of the \(H\) horizon-wise RMSE values and therefore differs from the overall test-set RMSE in Table~\ref{tab:benchmark-short}.

\subsection{Ablation Study}

The ablations ask whether the gain requires true pre-origin power, multi-scale encoding, and queryable weather memory, and whether weather remains dominant. Weather-only removes power context; Random power and Zero power retain the power path but destroy its state information, separating content from capacity. Basic GCN replaces multi-scale encoding, while the remaining variants isolate weather-memory organization and meteorological groups.

Table~\ref{tab:ablation-short} reports the results, where \(\Delta\) is computed from the unrounded 12-run means. Weather-only increases MAE and RMSE by 7.34\% and 7.04\%, respectively. Random power and Zero power degrade by similar amounts, showing that the benefit requires the numerical content of the true operating state rather than extra parameters or a power-side path alone. Basic GCN increases the two errors by 4.88\% and 4.74\%, indicating that multi-scale encoding strengthens the state representation but is secondary to providing a valid state in the first place.

\begin{table}[htbp]
\caption{Key ablation results}
\begin{center}
\footnotesize
\setlength{\tabcolsep}{2pt}
\begin{tabular}{lccccc}
\hline
Variant & MAE \(\downarrow\) & \(\Delta\)MAE & RMSE \(\downarrow\) & \(\Delta\)RMSE & \(R^2\) \(\uparrow\) \\
\hline
Ours & \textbf{0.0779} & 0.00\% & \textbf{0.1065} & 0.00\% & \textbf{0.8036} \\
Random power & 0.0835 & 7.17\% & 0.1137 & 6.78\% & 0.7760 \\
Zero power & 0.0838 & 7.61\% & 0.1141 & 7.21\% & 0.7742 \\
Basic GCN & 0.0817 & 4.88\% & 0.1115 & 4.74\% & 0.7845 \\
Weather-only & 0.0836 & 7.34\% & 0.1140 & 7.04\% & 0.7749 \\
No memory & 0.0824 & 5.86\% & 0.1128 & 5.93\% & 0.7796 \\
No delta & 0.0810 & 3.96\% & 0.1110 & 4.24\% & 0.7865 \\
Single branch & 0.0843 & 8.22\% & 0.1162 & 9.12\% & 0.7660 \\
No synoptic & 0.0858 & 10.14\% & 0.1166 & 9.50\% & 0.7645 \\
No wind & 0.1485 & 90.68\% & 0.1991 & 87.02\% & 0.3123 \\
\hline
\end{tabular}
\label{tab:ablation-short}
\end{center}
\end{table}

The weather-side ablations No delta, No memory, Single branch, and No synoptic all cause stable degradation, showing that meteorological changes, spatiotemporal organization, variable-role separation, and background fields contribute. No wind increases MAE and RMSE by 90.68\% and 87.02\%, respectively, and reduces \(R^2\) to 0.3123, confirming that wind fields remain the primary driver rather than the power context replacing future weather.

Together with Table~\ref{tab:benchmark-short}, these controls isolate the main choices. True/random/zero-power comparisons establish the need for valid state information; output-level fusion supports conditioning memory rather than merely adding a power predictor; Basic GCN quantifies the extra benefit of multi-scale encoding; and weather ablations confirm that the memory remains weather-driven. The evidence therefore supports the proposed information organization rather than component accumulation.

\section{Discussion and Conclusion}

Three-year experiments show that future weather remains the dominant intraday driver, while true pre-origin power provides an effective operating-state constraint. Conditioning weather-memory reading on that state improves average accuracy and substantially reduces initial-horizon error. The results further demonstrate how explicit information-role assignment and memory-level conditioning coordinate the forecast-origin state with future-weather information. Complementary baselines and controls support the mechanism and the secondary benefit of multi-scale encoding. Since the current evidence covers deterministic forecasting in one region, future work will examine weather uncertainty, cross-region transfer, model simplification, and extreme ramps.

\begin{thebibliography}{29}

\bibitem{ref1}
H. Liu, X. Mi, and Y. Li, ``Smart deep learning based wind speed prediction model using wavelet packet decomposition, convolutional neural network and convolutional long short term memory network,'' \textit{Energy Convers. Manag.}, vol. 166, pp. 120--131, 2018.

\bibitem{ref2}
R. Yu, Z. Liu, X. Li, et al., ``Scene learning: Deep convolutional networks for wind power prediction by embedding turbines into grid space,'' \textit{Appl. Energy}, vol. 238, pp. 249--257, 2019.

\bibitem{ref3}
L. Wang, J. Gao, Y. Li, and D. Wang, ``A method for ultra-short-term wind power forecasting of large-scale wind farms based on adaptive spatiotemporal graph convolution,'' \textit{Renew. Energy}, vol. 249, Art. no. 123321, 2025.

\bibitem{ref4}
M. Li, M. Yang, and Y. Yu, ``Short-term regional wind power forecast method based on hierarchical graph neural network and extreme value theory,'' \textit{Energy Convers. Manag.}, vol. 341, Art. no. 120002, 2025.

\bibitem{ref5}
X. Liu, Y. Zhang, Z. Zhen, F. Xu, F. Wang, and Z. Mi, ``Spatio-temporal graph neural network and pattern prediction based ultra-short-term power forecasting of wind farm cluster,'' \textit{IEEE Trans. Ind. Appl.}, vol. 60, no. 1, pp. 1794--1803, 2024.

\bibitem{ref6}
M. Yang, C. Ju, Y. Huang, Y. Guo, and M. Jia, ``Short-term power forecasting of wind farm cluster based on global information adaptive perceptual graph convolution network,'' \textit{IEEE Trans. Sustain. Energy}, vol. 15, no. 3, pp. 2063--2076, 2024.

\bibitem{ref7}
M. Yang, D. Wang, W. Zhang, and X. Yv, ``A centralized power prediction method for large-scale wind power clusters based on dynamic graph neural network,'' \textit{Energy}, vol. 310, Art. no. 133210, 2024.

\bibitem{ref8}
M. Yang, Y. Guo, T. Huang, et al., ``Wind farm cluster power prediction based on graph deviation attention network with learnable graph structure and dynamic error correction during load peak and valley periods,'' \textit{Energy}, vol. 312, Art. no. 133645, 2024.

\bibitem{ref9}
M. Yang, Y. Huang, Y. Guo, W. Zhang, and B. Wang, ``Ultra-short-term wind farm cluster power prediction based on FC-GCN and trend-aware switching mechanism,'' \textit{Energy}, vol. 290, Art. no. 130238, 2024.

\bibitem{ref10}
Z. Qu, J. Li, X. Hou, and J. Gui, ``A D-stacking dual-fusion, spatio-temporal graph deep neural network based on a multi-integrated overlay for short-term wind-farm cluster power multi-step prediction,'' \textit{Energy}, vol. 281, Art. no. 128289, 2023.

\bibitem{ref11}
M. Yang, Y. Jiang, C. Xu, et al., ``Day-ahead wind farm cluster power prediction based on trend categorization and spatial information integration model,'' \textit{Appl. Energy}, vol. 388, Art. no. 125580, 2025.

\bibitem{ref12}
Q. Zhang, K. Mei, Q. Zhou, et al., ``A wind farm cluster power prediction method based on globally aware dynamic spatiotemporal graph continual learning framework,'' \textit{Energy}, vol. 348, Art. no. 140591, 2026.

\bibitem{ref13}
L. Wang, Y. He, X. Liu, L. Li, and K. Shao, ``M2TNet: Multi-modal multi-task Transformer network for ultra-short-term wind power multi-step forecasting,'' \textit{Energy Rep.}, vol. 8, pp. 7628--7642, 2022.

\bibitem{ref14}
L. Wang and Y. He, ``M2STAN: Multi-modal multi-task spatiotemporal attention network for multi-location ultra-short-term wind power multi-step predictions,'' \textit{Appl. Energy}, vol. 324, Art. no. 119672, 2022.

\bibitem{ref15}
C. Liu, X. Zhang, S. Mei, Q. Zhou, and H. Fan, ``Series-wise attention network for wind power forecasting considering temporal lag of numerical weather prediction,'' \textit{Appl. Energy}, vol. 336, Art. no. 120815, 2023.

\bibitem{ref16}
M. Yang, Y. Guo, T. Huang, and W. Zhang, ``Power prediction considering NWP wind speed error tolerability: A strategy to improve the accuracy of short-term wind power prediction under wind speed offset scenarios,'' \textit{Appl. Energy}, vol. 377, Art. no. 124720, 2025.

\bibitem{ref_liu_bayesian_uncertainty}
L. Liu, J. Liu, Y. Ye, et al., ``Ultra-short-term wind power forecasting based on deep Bayesian model with uncertainty,'' \textit{Renew. Energy}, vol. 205, pp. 598--607, 2023.

\bibitem{ref17}
M. Yang, C. Han, W. Zhang, and B. Wang, ``A short-term power prediction method for wind farm cluster based on the fusion of multi-source spatiotemporal feature information,'' \textit{Energy}, vol. 294, Art. no. 130770, 2024.

\bibitem{ref18}
S. Hu, Y. Xiang, H. Zhang, et al., ``Hybrid forecasting method for wind power integrating spatial correlation and corrected numerical weather prediction,'' \textit{Appl. Energy}, vol. 293, Art. no. 116951, 2021.

\bibitem{ref19}
K. Higashiyama, Y. Fujimoto, and Y. Hayashi, ``Feature extraction of NWP data for wind power forecasting using 3D-convolutional neural networks,'' \textit{Energy Procedia}, vol. 155, pp. 350--358, 2018.

\bibitem{ref20}
J. Keisler and E. Le Naour, ``WindDragon: automated deep learning for regional wind power forecasting,'' \textit{Environ. Data Sci.}, vol. 4, Art. no. e19, 2025.

\bibitem{ref21}
J. Jonkers, D. N. Avendano, G. Van Wallendael, and S. Van Hoecke, ``A novel day-ahead regional and probabilistic wind power forecasting framework using deep CNNs and conformalized regression forests,'' \textit{Appl. Energy}, vol. 361, Art. no. 122900, 2024.

\bibitem{ref22}
K. Bi, L. Xie, H. Zhang, C. Chen, X. Gu, and Q. Tian, ``Accurate medium-range global weather forecasting with 3D neural networks,'' \textit{Nature}, vol. 619, no. 7970, pp. 533--538, 2023.

\bibitem{ref23}
H. Zhou, S. Zhang, J. Peng, et al., ``Informer: Beyond efficient Transformer for long sequence time-series forecasting,'' in \textit{Proc. AAAI Conf. Artif. Intell.}, vol. 35, no. 12, pp. 11106--11115, 2021.

\bibitem{ref24}
Y. Nie, N. H. Nguyen, P. Sinthong, and J. Kalagnanam, ``A time series is worth 64 words: Long-term forecasting with Transformers,'' in \textit{Proc. Int. Conf. Learn. Representations (ICLR)}, 2023.

\bibitem{ref25}
Y. Zhang and J. Yan, ``Crossformer: Transformer utilizing cross-dimension dependency for multivariate time series forecasting,'' in \textit{Proc. Int. Conf. Learn. Representations (ICLR)}, 2023.

\bibitem{ref26}
Y. Liu, H. Wu, J. Wang, and M. Long, ``Non-stationary Transformers: Exploring the stationarity in time series forecasting,'' in \textit{Adv. Neural Inf. Process. Syst.}, vol. 35, pp. 9881--9893, 2022.

\bibitem{ref_liu_feature_temporal}
L. Liu, X. Wang, X. Dong, et al., ``Interpretable feature-temporal transformer for short-term wind power forecasting with multivariate time series,'' \textit{Appl. Energy}, vol. 374, Art. no. 124035, 2024.

\bibitem{ref_liu_tdformer}
L. Liu, X. Liu, J. Huang, et al., ``TDFormer: A novel triple decoupled transformer for accurate multi-step wind power forecasting,'' \textit{Eng. Appl. Artif. Intell.}, vol. 181, Art. no. 115310, 2026.

\end{thebibliography}


\end{document}
